\chapter{An anisotropic band-limit}
\label{ch:method1}

\keybox{the idea}{%
Think of each blob's position as something the photographs \emph{measure}. Every training photograph that sees the blob is a
measurement of where it is, and every measurement has noise --- roughly a pixel's
worth, in the image plane.

Now ask the standard estimation question: given all those measurements, how precisely do you know the position? The answer varies with direction. A camera
localises a point very well within its image plane and poorly along its own
viewing direction; a second camera at a different angle fixes the first camera's
weak direction. Accumulate over all the cameras that saw the blob and you get a
$3\times3$ matrix --- the familiar Fisher information of the measurement model, written $\Lambda_k$
throughout. Its eigenvectors are the directions the capture
constrains well and badly; its eigenvalues say how well.

The proposal is then almost forced: \emph{a primitive should not be finer, in any
direction, than the capture can determine it in that direction}. Invert the
matrix, and you have a covariance floor with the right shape. Then floor that at the per-view Nyquist limit, since no amount of accumulated
information lets you resolve structure finer than a pixel, and you have the
method.

Three things follow, and this chapter proves them. Mip-Splatting's scalar filter
is what you get by collapsing this matrix to a single view and throwing
away direction. Its mysterious constant $\sqrt{0.2}$ turns out to be a confidence
radius of $0.447$ pixels. And when the Nyquist floor dominates, which on a normal benchmark capture it
always does, the new method becomes the old one exactly. That last point is why the standard benchmarks in
Chapter~\ref{ch:results} show no improvement, and why a gain there would have indicated a bug.}

\begin{figure}[htbp]
  \centering
  \includegraphics[width=0.97\linewidth]{figures/d3-one-camera-two-cameras.pdf}
  \caption{Why the limit has to be a matrix. A single camera fixes a point well
  across its line of sight and badly along it, so the region consistent with its
  measurement is a long thin cigar pointing back at the lens --- the familiar
  fact that one photograph gives you direction but not distance. A second camera from elsewhere is precise where the first is vague, and the region
  consistent with both is compact. Accumulating this over every camera that saw
  the primitive is what $\Lambda_k$ does, and its eigenvectors are these axes.
  A scalar limit cannot distinguish the long axis from the short one; that is
  the whole of the argument.}
  \label{fig:cameras}
\end{figure}

\section{The estimation problem}
Regard the position $p_k$ of a primitive as a parameter to be estimated from
image measurements. Camera $n$ observes the primitive's projection at
\begin{equation}
  u_n \;=\; \pi_n(p_k) + \varepsilon_n,\qquad
  \varepsilon_n \sim \mathcal{N}(0,\Sigma_{\mathrm{pix}}),\quad
  \Sigma_{\mathrm{pix}} = \sigma_{\mathrm{pix}}^2 I_2 ,
  \label{eq:measurement}
\end{equation}
where $\sigma_{\mathrm{pix}}$ is the localisation noise of an image feature, in
pixels. The log-likelihood of the measurement from view $n$ is, up to constants,
$-\tfrac12 (u_n-\pi_n(p))^{\!\top}\Sigma_{\mathrm{pix}}^{-1}(u_n-\pi_n(p))$.

\section{The measurement Jacobian}
\label{sec:jacobian}
The Fisher information carried by view $n$ about $p$ is
\begin{equation}
  I_n \;=\; \Big(\frac{\partial \pi_n}{\partial p}\Big)^{\!\top}
            \Sigma_{\mathrm{pix}}^{-1}
            \Big(\frac{\partial \pi_n}{\partial p}\Big),
  \qquad \frac{\partial \pi_n}{\partial p} = J_n W_n .
  \label{eq:fisher-view}
\end{equation}
The derivative of image position with respect to world position is the EWA pair
of Equation~\eqref{eq:ewa}. \emph{No new geometry is introduced}: the quantity
the renderer forms to project a covariance is the same quantity estimation
theory requires to bound a position. This identification is the technical core
of the thesis, and it is what makes the method cheap.

\paragraph{Rank.} $J_nW_n \in \mathbb{R}^{2\times3}$, so $I_n$ has rank at most
$2$. A single view constrains the two directions transverse to the ray and says
nothing about depth. For a pinhole camera the eigenvalues of $I_n$ are exactly
$(0,\ (f/z)^2/\sigma_{\mathrm{pix}}^2,\ (f/z)^2/\sigma_{\mathrm{pix}}^2)$
(Appendix~\ref{app:rank}). Depth information exists only through parallax.

\section{The observability matrix}
\label{sec:accumulation}
Summing over views, weighted by each view's visibility and blending contribution $w_{nk}$, which the rasteriser already accumulates,
\begin{equation}
  \Lambda_k \;=\; \sum_n w_{nk}\,(J_nW_n)^{\!\top}\Sigma_{\mathrm{pix}}^{-1}(J_nW_n)
  \;=\; V\,\mathrm{diag}(\lambda_1,\lambda_2,\lambda_3)\,V^{\!\top}.
  \label{eq:lambda}
\end{equation}

\paragraph{The weight, defined.} $w_{nk} = T_{nk}\,\alpha_k \in [0,1]$ is the
primitive's blending contribution in view $n$: the transmittance accumulated in
front of it times its own opacity, which is the coefficient the rasteriser already multiplies its colour by. It is zero for a primitive the view
does not see and near one for an unoccluded opaque one. This weighting is a \emph{heuristic} down-weighting of occluded views, standing outside any
likelihood:
a view that cannot see a primitive carries no information about it, and $w_{nk}$
is the cheapest available statement of that fact. It is stated as a heuristic
because calling it anything more would be a claim the derivation does not
support.

\paragraph{What the matrix is, precisely.} With unit weights,
Equation~\eqref{eq:lambda} \emph{is} the Fisher information of the measurement
model in Equation~\eqref{eq:measurement}. With the weights it departs from that, so the honest name for $\Lambda_k$ is an
\textbf{observability matrix} --- a
view-weighted accumulation of the per-view measurement information, coinciding
with the Fisher information exactly when $w_{nk} \equiv 1$. Under the
point-feature approximation of Equation~\eqref{eq:measurement},
Equation~\eqref{eq:lambda} is also the Gauss--Newton approximation to the Hessian
of the photometric cost in $p_k$: the residual is linearised through $J_nW_n$ and second-derivative terms are dropped, as every bundle adjustment does and for
the same reason.

\paragraph{Why this matters for the claim.} The Cram\'er--Rao bound
$\mathrm{Cov}(\hat p_k) \succeq \Lambda_k^{-1}$ holds for an unbiased estimator
of the measurement model as written. A splatting optimiser is not that estimator:
it fits a photometric loss over a densifying set of correlated primitives, not
independent point features with Gaussian localisation noise. Cram\'er--Rao is
therefore kept as the \emph{motivation} for the form of Equation~\eqref{eq:filtervar}: it says why
the inverse of an information matrix is the right shape for a resolution limit.
The claim the experiments actually test is weaker, and sufficient: along eigendirection $i$ the capture constrains position
to order $1/\sqrt{\lambda_i}$, so structure finer than that is not identifiable
from these images.

\keybox{principle}{%
A primitive should not be narrower, along a given direction, than the
uncertainty with which the capture determines its position along that direction.
Structure finer than $s/\sqrt{\lambda_i}$ is not identifiable from the
measurements; fitting it is fitting noise.}

\section{The covariance floor}
Hence a per-direction variance, floored at Mip-Splatting's own value and capped
at a scene-scale fraction:
\begin{equation}
  \sigma_i^2 \;=\; \min\!\Big(\max\big(f_k^2,\ s^2/\lambda_i\big),\ (\beta z_{\min,k})^2\Big),
  \label{eq:filtervar}
\end{equation}
\begin{equation}
  \Sigma_{\mathrm{filt}} = V\,\mathrm{diag}(\sigma_1^2,\sigma_2^2,\sigma_3^2)\,V^{\!\top};
  \quad \Sigma_k \leftarrow \Sigma_k + \Sigma_{\mathrm{filt}};
  \quad \alpha_k \leftarrow \alpha_k\sqrt{|\Sigma_k|/|\Sigma_k+\Sigma_{\mathrm{filt}}|}.
  \label{eq:filter}
\end{equation}
Two hyperparameters are introduced and none removed. $\beta$ caps the filter in
degenerate directions where $\lambda_i \to 0$; without it an unconstrained
direction would produce an unbounded primitive. The opacity compensation is
Mip-Splatting's own, unchanged, and Appendix~\ref{app:mass} shows it carries over
to an anisotropic $\Sigma_{\mathrm{filt}}$ without modification.

\paragraph{$s$ and $\sigma_{\mathrm{pix}}$, disentangled.} These two appear only
as the product $s\,\sigma_{\mathrm{pix}}$, so quoting a default for one and a
value for the product over-determines them. The convention used throughout, and
the one the implementation follows:
\begin{itemize}[leftmargin=1.4em,itemsep=2pt]
  \item $\sigma_{\mathrm{pix}} = \sqrt{0.2} = 0.447$\,px is \emph{inherited} from Mip-Splatting's own constant. Proposition~\ref{prop:constant} shows it is the localisation
        noise implicit in Mip-Splatting's constant, and taking it verbatim is
        what makes Proposition~\ref{prop:reduction} exact instead of merely close.
  \item $s$ is the dimensionless confidence multiplier, default $\mathbf{1}$
        (one standard deviation). It is the only free knob of the pair, and
        every result in this thesis is at $s=1$.
\end{itemize}
So $s\,\sigma_{\mathrm{pix}} = 0.447$\,px at the default, and a reader changing
$s$ changes the confidence level and nothing else.

\paragraph{Consistency of the floor and the cap.} The floor must never exceed the
cap, or Equation~\eqref{eq:filtervar} would be ill-posed. Requiring
$f_k \le \beta z_{\min,k}$ and substituting Equation~\eqref{eq:mipfilter} gives
$\beta \ge \sqrt{0.2}/f_{\max}$. For $f_{\max}\approx1000$\,px this is
$\beta \ge 4.5\times10^{-4}$; the default $\beta=0.01$ carries $22\times$
headroom.

\section{Mip-Splatting as a special case}
\label{sec:proposition1}
Take a single camera with focal length $f$ viewing a primitive at depth $z$ on
the optical axis. From \S\ref{sec:jacobian} the lateral eigenvalue is
$\lambda_{\mathrm{lat}} = (f/z)^2/\sigma_{\mathrm{pix}}^2$, so the
estimation-limited lateral scale is $s/\sqrt{\lambda_{\mathrm{lat}}} =
s\,\sigma_{\mathrm{pix}}\,z/f$.

\begin{proposition}[Identification of the constant]
\label{prop:constant}
Mip-Splatting's 3D smoothing filter is exactly Equation~\eqref{eq:filtervar}
evaluated for a single view, restricted to the lateral directions, and applied
isotropically --- provided
\[ s\cdot\sigma_{\mathrm{pix}} \;=\; \sqrt{0.2} \;\approx\; 0.447\ \text{pixels}. \]
\end{proposition}

Their constant therefore carries a meaning: a sub-pixel localisation confidence
of $0.447$\,px, slightly under half a pixel. The present method keeps
that calibration and generalises the geometry around it. This is the strongest
single argument in the thesis: it converts the relationship between the proposed
method and its baseline from ``similar in spirit'' into ``the baseline is a
corollary'', and it explains a magic number in a published paper.
Appendix~\ref{app:prop1} gives the derivation in full.

\section{Off-axis anisotropy}
\label{sec:offaxis}

Before parallax enters at all, there is anisotropy the scalar filter discards.
Section~\ref{sec:jacobian} quotes the single-view eigenvalues
$(0, (f/z)^2, (f/z)^2)/\sigma_{\mathrm{pix}}^2$, and that is exact --- \emph{on
the optical axis}. The convenient rank-1 form
$(f/z)^2(I - dd^{\!\top})/\sigma_{\mathrm{pix}}^2$ suggests otherwise, because
$I - dd^{\!\top}$ is a projector and a projector's non-zero eigenvalues are equal
by construction, whatever $d$ is. Off-axis, the true $J_n^{\!\top}J_n$ departs from a projector, and the two
lateral eigenvalues separate.

Writing $a = (x/z,\,y/z)$ for the normalised off-axis displacement,
$J = (f/z)B$ with $B = \big[\begin{smallmatrix}1&0&-a_1\\0&1&-a_2\end{smallmatrix}\big]$,
so $BB^{\!\top}$ has trace $2+|a|^2$ and determinant $1+|a|^2$. Its eigenvalues
are therefore $1+|a|^2$ and $1$, giving
\begin{equation}
  \frac{\lambda_2}{\lambda_1} \;=\; \frac{1}{1+|a|^2} \;=\; \cos^2\varphi,
  \qquad \varphi = \arctan|a| \ \text{the off-axis angle.}
  \label{eq:offaxis}
\end{equation}

\input{generated/table-geometry}

At the corner of a $50^\circ$-field camera $\varphi$ approaches $25^\circ$ and
the two lateral eigenvalues already differ by $20\,\%$; the sampled poses of
Table~\ref{tab:offaxis} run from $0.93$ down to $0.47$. So \textbf{perspective
contributes anisotropy independently of parallax}, and a scalar filter discards
it before the parallax argument of \S\ref{sec:anisotropy} has even begun. The
implementation uses the full $J_nW_n$ and never the rank-1 shortcut, so this term survives into the filter, where the shortcut would discard it for
good; the shortcut appears in this document
only where it makes an argument shorter, and Equation~\eqref{eq:offaxis} records
what it costs. This is also the most likely explanation for the residual in
\S\ref{sec:instrumentresults}: the grazing protocol, which is a full orbit tilted
towards obliquity, measures more anisotropy than any parallax-only model
predicts.

\section{The two-view spectrum}
\label{sec:anisotropy}
Take the simplest case: two cameras at the same distance $z$, placed symmetrically about the primitive, so that the angle
between their lines of sight is $\theta$. Evaluating
Equation~\eqref{eq:lambda} for this arrangement (the algebra is in
Appendix~\ref{app:twoview}) gives $\Lambda \propto 2I - d_1d_1^{\!\top} -
d_2d_2^{\!\top}$, and this has the exact spectrum
$\{2,\ 2\cos^2(\theta/2),\ 2\sin^2(\theta/2)\}$.

These are three distinct numbers. The right filter is therefore triaxial, with a different limit along each of
three axes: the lateral
direction lying in the plane of the two rays is degraded as well, and only the
direction perpendicular to that plane is fully constrained. Taking the extreme
pair of eigenvalues,
\begin{equation}
  \frac{\lambda_3}{\lambda_1} = \sin^2(\theta/2),
  \qquad
  \frac{\sigma_{\mathrm{depth}}}{\sigma_{\mathrm{lat}}} = \frac{1}{\sin(\theta/2)}
  \qquad (\theta \le 90^\circ).
  \label{eq:anisotropy}
\end{equation}
This is exact at every angle --- verified against pinhole Jacobians to
$4\times10^{-16}$ over $\theta = 2^\circ$ to $179^\circ$
(Table~\ref{tab:twoviewspectrum}). The restriction $\theta \le 90^\circ$ earns its place: past $90^\circ$ the depth direction becomes better constrained than the
in-plane lateral one and the two exchange roles, so the ratio of smallest to
largest is $\min(\sin^2(\theta/2), \cos^2(\theta/2))$. At $\theta = 120^\circ$
the naive form reads $0.75$ where the truth is $0.25$ --- a factor of three, in
the optimistic direction. Every protocol this thesis stresses lies well below
$90^\circ$, and the full orbit lies so far above it that the ratio is $\approx 1$
either way, so the restriction costs nothing here and guards against a real error elsewhere.

\subsection{The spherical-cap correction}
\label{sec:capprediction}

Equation~\eqref{eq:anisotropy} is exact for \emph{two} cameras. A capture
protocol is ten or a hundred, spread over an angular window, and $\Lambda$ is the
sum over that spread --- a different number. Reading the two-view formula off a protocol's measured span is defensible only
for a protocol that really is two cameras, and every protocol here has more.

For cameras distributed uniformly over a spherical cap of half-angle
$\alpha = \theta/2$ about the mean viewing direction, $\cos t$ is uniform on
$[\cos\alpha, 1]$, so with $c = \cos\alpha$ and $E = \mathbb{E}[d_z^2] =
(1+c+c^2)/3$,
\begin{equation}
  \frac{\lambda_3}{\lambda_1} \;=\; \frac{2(1-E)}{1+E},
  \qquad
  \frac{\sigma_{\mathrm{depth}}}{\sigma_{\mathrm{lat}}}
  \;=\; \sqrt{\frac{1+E}{2(1-E)}} .
  \label{eq:cap}
\end{equation}

\keybox{the two-view form against a real capture}{%
As $\alpha \to 0$, Equation~\eqref{eq:cap} tends to $\alpha^2/2$ while
$\sin^2\alpha \to \alpha^2$: the cap's $\lambda_3/\lambda_1$ is
\textbf{half} the two-view value, because a cap's cameras are spread across the window instead of sitting at its two ends. A smaller $\lambda_3/\lambda_1$ is a
\emph{worse}-conditioned capture, so Equation~\eqref{eq:anisotropy} applied to a
protocol \emph{understates} its depth anisotropy by $\sqrt{2}$ in $\sigma$. The
direction matters: the conditioning instrument measuring more anisotropy than the two-view form predicts
is the model working, not failing, and reporting it against the two-view number
would have made a confirmation look like a discrepancy. Every prediction quoted
against a measured span in this thesis therefore uses Equation~\eqref{eq:cap}.}

\begin{figure}[htbp]
  \centering
  \includegraphics[width=\linewidth]{figures/fig5-fisher-geometry.pdf}
  \caption{The mechanism. Cross-sections of the two filters in the plane of the
  cameras, normalised to the same lateral extent because the floor makes them
  agree there by construction, so the whole visible difference is along depth.
  The two panels are the spans the \texttt{cone} and \texttt{arc} protocols
  actually realise ($20^\circ$ and $80^\circ$), and the elongation is
  Equation~\eqref{eq:cap} at those spans --- $8.1\times$ and $2.0\times$. On a
  full orbit the two ellipses coincide to within the line width, which is Proposition~\ref{prop:reduction} drawn instead of stated. Geometry computed from real pinhole Jacobians.}
  \label{fig:fisher}
\end{figure}

\begin{figure}[htbp]
  \centering
  \includegraphics[width=0.86\linewidth]{figures/fig6-anisotropy.pdf}
  \caption{Where the method has room, and which prediction to use. Mip-Splatting
  assumes the ratio is $1$ at every span; the gap between that line and the
  curve is the quantity this thesis exploits. Both forms are drawn because the
  difference between them is a correction this project had to make: the dashed
  curve is Equation~\eqref{eq:anisotropy}, exact for two cameras, and the solid
  one is Equation~\eqref{eq:cap} for cameras distributed over the span, as a protocol is. The points are the conditioning instrument measured on eight scenes
  (\S\ref{sec:instrumentresults}), and they sit above the two-view curve at
  every protocol --- plotted against that curve alone, a confirmation would read
  as a discrepancy.}
  \label{fig:anisotropy}
\end{figure}

\begin{table}[htbp]
  \centering
  \caption{What Equation~\eqref{eq:cap} predicts, at the spans the protocols of
  \S\ref{sec:stress} actually realise. A set of falsifiable predictions, made
  before the experiments, and the reason the stress suite exists. A method that helped everywhere would contradict its own derivation.}
  \label{tab:regimes}
  \small
  \begin{tabular}{lccl}
    \toprule
    capture & span & $\sigma_D/\sigma_L$ & expectation \\
    \midrule
    full orbit (standard benchmarks) & $179^\circ$ & $1.00\times$ & parity --- the floor dominates \\
    grazing orbit & $168^\circ$ & $1.04\times$ & parity from parallax alone \\
    one-sided arc & $80^\circ$ & $2.03\times$ & gain along depth \\
    low-parallax cone & $20^\circ$ & $8.10\times$ & largest gain \\
    single view & $0^\circ$ & $\infty$ & cap $\beta$ binds; scene-scale prior \\
    \bottomrule
  \end{tabular}
  \par\vspace{2pt}\footnotesize Spans are the medians the built protocols
  measured (Table~\ref{tab:instruments}), so the prediction and the measurement
  are evaluated at the same geometry. Quoting Equation~\eqref{eq:anisotropy}
  here instead would give $1.00$, $1.01$, $1.56$ and $5.76\times$ --- uniformly
  low, for the reason \S\ref{sec:capprediction} gives.
\end{table}

\section{Exact reduction}
\label{sec:proposition2}

\begin{figure}[htbp]
  \centering
  \includegraphics[width=0.97\linewidth]{figures/d5-floor-vs-estimate.pdf}
  \caption{Why a benchmark result of ``no change'' is the correct outcome. Two limits are in play and the method takes whichever
  is larger. The flat blue line is the resolution floor: no capture, however
  generous, lets you resolve structure finer than a pixel. The rising orange
  curve is the estimation limit of this chapter, which grows as the cameras
  crowd into a narrower window. On a standard benchmark, meaning a hundred cameras around a full circle, the
  floor sits far above the estimation limit, the
  method returns the floor, and it \emph{is} Mip-Splatting, exactly and provably
  (Proposition~\ref{prop:reduction}). Only when the capture narrows past the
  crossing does anything change. The four marks are the protocols of
  \S\ref{sec:stress} at their measured widths; the crossing falls between the
  cone and the pencil, which is where the measurement of \S\ref{sec:table3} puts
  it.}
  \label{fig:floorvsest}
\end{figure}

\begin{proposition}[Exact reduction]
\label{prop:reduction}
If $s^2/\lambda_i \le f_k^2$ for every $i$, then Equation~\eqref{eq:filter} is
identical to Mip-Splatting's Equation~\eqref{eq:mipfilter}.
\end{proposition}

\begin{proof}
Under the hypothesis, and given $\beta \ge \sqrt{0.2}/f_{\max}$ so the cap does
not bind, Equation~\eqref{eq:filtervar} gives $\sigma_i^2 = f_k^2$ for all $i$.
Then $\Sigma_{\mathrm{filt}} = V(f_k^2 I)V^{\!\top} = f_k^2 VV^{\!\top} = f_k^2 I$,
since $V$ is orthogonal. The opacity compensation of Equation~\eqref{eq:filter}
then coincides with Equation~\eqref{eq:opacitycomp}.
\end{proof}

\paragraph{Consequence.} The filter never smooths \emph{less} than the baseline
in any direction. Parity on well-conditioned benchmarks is therefore a
structural property, not an empirical hope, and a failure to achieve it is a
bug, testable to numerical precision.

This proposition is implemented as a unit test: with the accumulation disabled the model must reproduce Mip-Splatting to $10^{-6}$,
and with the filter zeroed entirely it must reproduce 3DGS. The second of these
is already in force. It \emph{is} the 3DGS-arm construction of \S\ref{sec:onecodebase}, and Chapter~\ref{ch:results} reports it verified on the saved models.

\section{Implementation cost}
\label{sec:implcost}

Two consequences follow from $\Sigma_{\mathrm{filt}}$ being a matrix. Both belong to the implementation.

\paragraph{The full $3\times3$ must reach the rasteriser.} Mip-Splatting can add
its filter by updating each scale as $s_i \leftarrow \sqrt{s_i^2 + f_k^2}$,
because $f_k^2 I$ commutes with $\Sigma_k$ and shares its eigenvectors. An
anisotropic $\Sigma_{\mathrm{filt}}$ does not commute with $\Sigma_k$ in general:
$V$ diagonalises $\Lambda_k$ and $R$ diagonalises $\Sigma_k$, and nothing aligns
them. So $\Sigma_k + \Sigma_{\mathrm{filt}}$ is not expressible as a per-axis
update of the scale vector, and the summed covariance has to be formed in full
and passed through the rasteriser's precomputed-covariance input. The implementation does that. Where the filter comes out isotropic, which Proposition~\ref{prop:reduction}
makes the common case, it takes Mip-Splatting's own scalar path instead, so parity holds in the implementation as well as in the algebra: the two take
the same code path.

\paragraph{The eigendecomposition stays off the backward path.} $\Lambda_k$ is
computed under \texttt{no\_grad} and held constant between recomputations,
exactly as \texttt{filter\_3D} is. The gradient therefore never passes through
it, and the question of differentiating a possibly degenerate eigenvector basis
never arises, though it is where a naive implementation of
this idea would have failed. The filter is a piecewise-constant
band-limit that the optimiser sees as a fixed additive covariance, held constant throughout
training.

\section{Arithmetic cost}
$\Lambda_k$ is a sum of $N$ outer products of a $2\times3$ matrix, followed by a
$3\times3$ symmetric eigendecomposition per primitive. It is recomputed on the
same schedule Mip-Splatting already uses for \texttt{filter\_3D} --- every 100
iterations, roughly 290 times across a 30\,000-iteration run, and never at
render time. Rendering speed is therefore unchanged by construction, which is
the answer to the first question a graphics audience will ask, and \S\ref{sec:fps} measures it instead of leaving it as an argument.
